\documentclass[aspectratio=169]{beamer}
\usetheme{simple}

\input{preamble/preamble.tex}
\input{preamble/preamble_math.tex}
% \input{preamble/preamble_acronyms.tex}

\title{Lecture 7: Markov Chain Monte Carlo}
\subtitle{STATS305B: Applied  Statistics II}
\author{Scott Linderman}
\date{\today}


\begin{document}


\maketitle

\begin{frame}{Announcements}

    \begin{itemize}
        \item HW2 posted online. Due Wed, Feb 12.
        \item \textcolor{red}{\textbf{Midterm next Wednesday in MCCULL 115}} See practice test online.
    \end{itemize}
    
\end{frame}

\begin{frame}{Recap}

    Last Time...
    \begin{itemize}
        \item Bayesian Inference
        \item Conjugate Priors
        \item Laplace Approximation
        \item Monte Carlo
        \item Start on MCMC
    \end{itemize}
    
\end{frame}

\begin{frame}{Outline}

  Today...
  \begin{itemize}
      \item Markov Chains
      \item Metropolis Hastings Algorithm
      \item Gradient-based Proposals (MALA and HMC)
      \item Gibbs Sampling
      \item Augmentation Schemes (with Demo)
  \end{itemize}
  
\end{frame}
    

    
\begin{frame}[t,allowframebreaks]{Approximating Posterior
Expectations}\label{approximating-posterior-expectations}

Generally, we can't analytically compute posterior expectations. In
these cases, we need to resort to approximations. For example, we could
use \emph{quadrature methods} like Simpson's rule or the trapezoid rule
to numerically approximate the integral over \(\Theta\).

Roughly, \begin{align*}
\E_{p(\theta | x)}[f(\theta)] \approx \sum_{m=1}^M p(\theta_m \mid x) \, f(\theta_m) \, \Delta_m
\end{align*} where \(\theta_m \subset \Theta\) is a grid of points and
\(\Delta_m\) is a volume around that point.

This works for low-dimensional problems (say, up to \(5\) dimensions),
but the number of points (\(M\)) needed to get a good estimate grows
exponentially with the parameter dimension.
\end{frame}

\begin{frame}[t,allowframebreaks]{Monte Carlo
Approximations}\label{monte-carlo-approximations}

\textbf{Idea:} approximate the expectation via sampling,

\begin{align*}
\E_{p(\theta | x)}[f(\theta)] \approx \frac{1}{M} \sum_{m=1}^M f(\theta_m) \quad \text{where} \quad \theta_m \sim p(\theta \mid x).
\end{align*}

Let \(\hat{f} = \frac{1}{M} \sum_{m=1}^M f(\theta_m)\) denote the Monte
Carlo estimate. It is a random variable, since it's a function of random
samples \(\theta_m\). As such, we can reason about its mean and
variance.
\end{frame}

\begin{frame}[t,allowframebreaks]{Unbiasedness}\label{unbiasedness}

Clearly,

\begin{align*}
\E[\hat{f}] = \frac{1}{M} \sum_{m=1}^M \E_{p(\theta | x)}[f(\theta)] = \E_{p(\theta | x)}[f(\theta)].
\end{align*}

Thus, \(\hat{f}\) is an \emph{unbiased} estimate of the desired
expectation.
\end{frame}


\begin{frame}[t,allowframebreaks]{Monte Carlo Variance}\label{monte-carlo-variance}

What about its variance? \begin{align*}
\Var[\hat{f}] = \Var \left(\frac{1}{M} \sum_{m=1}^M f(\theta_m) \right) = \frac{1}{M^2} \left( \sum_{m=1}^M \Var[f(\theta)] + 2 \sum_{1 \leq m < m' \leq M} \mathrm{Cov} [f(\theta_m), f(\theta_{m'})] \right)
\end{align*}
\end{frame}

\begin{frame}[t,allowframebreaks]{Comparison to Numerical
Quadrature}\label{comparison-to-numerical-quadrature}

\begin{itemize}
\item
  If the samples are not only identically distributed but also
  \emph{uncorrelated}, then
  \(\Var[\hat{f}] = \frac{1}{M} \Var[f(\theta)]\).
\item
  In this case, the \emph{root mean squared error} (RMSE) of the
  estimate is \(\sqrt{\Var[\hat{f}]} = O(M^{-\frac{1}{2}})\).
\item
  Compare this to Simpson's rule, which for smooth 1D problems has an
  error rate of \(O(M^{-4})\). That's roughly 8 times better than Monte
  Carlo!
\item
  However, for multidimensional problems, Simpson's rule is
  \(O(M^{-\frac{4}{D}})\), whereas the \textbf{error rate of Monte Carlo
  does not depend on the dimensionality!}
\end{itemize}
\end{frame}

\begin{frame}[t,allowframebreaks]{The Catch}\label{the-catch}

So far so good: we'll just draw a lot of samples to drive down our Monte
Carlo error. \textbf{Here's the catch!} How do you draw samples from the
posterior \(p(\theta \mid x)\)? We're interested in Monte Carlo for
cases where the posterior does not admit a simple closed form! In
general, sampling the posterior is as hard as computing the marginal
likelihood.
\end{frame}

\begin{frame}[t,allowframebreaks]{Markov Chains}\label{markov-chains}

A \emph{Markov chain} is a joint distribution of a sequence of
variables, \(\pi(\theta_1, \theta_2, \ldots, \theta_M)\). (To avoid
confusion with the model \(p\), we denote the densities associated with
the Markov chain by \(\pi\).) The Markov chain factorizes so that each
variable is drawn conditional on the previous variable, \begin{align*}
\pi(\theta_1, \theta_2, \ldots, \theta_M) = \pi_{1}(\theta_1) \prod_{m=2}^M \pi(\theta_m \mid \theta_{m-1}).
\end{align*} This is called the \emph{Markov property}.

\begin{itemize}
\item
  The distribution \(\pi_1(\theta_1)\) is called the \emph{initial
  distribution}.
\item
  The distribution \(\pi(\theta_m \mid \theta_{m-1})\) is called the
  \emph{transition distribution}. If the transition distribution is the
  same for each \(m\), the Markov chain is \emph{homogeneous}.
\end{itemize}
\end{frame}

\begin{frame}[t,allowframebreaks]{Stationary distributions}\label{stationary-distributions}

Let \(\pi_m(\theta_m)\) denote the marginal distribution of sample
\(\theta_m\). It can be obtained recursively as, \begin{align*}
\pi_m(\theta_m) = \int \pi_{m-1}(\theta_{m-1}) \, \pi(\theta_m \mid \theta_{m-1}) \dif \theta_{m-1}.
\end{align*} We are interested in the asymptotic behavior of the
marginal distributions as \(m \to \infty\).

A distribution \(\pi^\star(\theta)\) is a \textbf{stationary
distribution} if, \begin{align*}
\pi^\star(\theta) = \int \pi^\star(\theta') \, \pi(\theta \mid \theta') \dif \theta'.
\end{align*} That is, suppose the marginal of sample \(\theta'\) is
\(\pi^\star(\theta)\). Then the marginal of the next time point is also
\(\pi^\star(\theta)\).
\end{frame}

\begin{frame}[t,allowframebreaks]{Detailed balance}\label{detailed-balance}

How can we relate transition distributions and stationary distributions?
A sufficient (but not necessary) condition for \(\pi^\star(\theta)\) to
be a stationary distribution is that it satisfies \emph{detailed
balance}, \begin{align*}
\pi^\star(\theta') \pi(\theta \mid \theta') = \pi^\star(\theta) \pi(\theta' \mid \theta).
\end{align*} In words, the probability of starting at \(\theta'\) and
moving to \(\theta\) is the same as that of starting at \(\theta\) and
moving to \(\theta'\), if you draw the starting point from the
stationary distribution.

To see that detailed balance is sufficient, integrate both sides to get,
\begin{align*}
\int \pi^\star(\theta') \pi(\theta \mid \theta') \dif \theta' = \int \pi^\star(\theta) \pi(\theta' \mid \theta) \dif \theta' = \pi^\star(\theta).
\end{align*} Thus, \(\pi^\star(\theta)\) is a stationary distribution of
the Markov chain with transitions \(\pi(\theta \mid \theta')\).
\end{frame}

\begin{frame}[t,allowframebreaks]{Ergodicity}\label{ergodicity}

Detailed balance can be used to show that \(\pi^\star(\theta)\) is
\emph{a} stationary distribution, but not that it is \emph{the unique}
one. This is where \emph{ergodicity} comes in. A Markov chain is ergodic
if \(\pi_m(\theta_m) \to \pi^\star(\theta)\) regardless of
\(\pi_1(\theta_1)\). An ergodic chain has only one stationary
distribution, \(\pi^\star(\theta)\).

The easiest way to prove ergodicity is to show that it is possible to
reach any \(\theta'\) from any other \(\theta\). E.g. this is trivially
so if \(\pi(\theta' \mid \theta) > 0\).

\textit{Note:} A more technical definition is that all pairs of
sets \emph{communicate}, in which case the chain is \emph{irreducible},
and that each state is \emph{aperiodic}. The definitions can be a bit
overwhelming. 
\end{frame}


\begin{frame}[t,allowframebreaks]{Markov Chain Monte Carlo
(MCMC)}\label{markov-chain-monte-carlo-mcmc}

Finally, we come to our \textbf{main objective}: designing a Markov
chain for which \emph{the posterior is the unique stationary
distribution.} That is, we want
\(\pi^\star(\theta) = p(\theta \mid x)\).

Recall our \textbf{constraint}: we can only compute the joint
probability (the numerator in Bayes' rule), not the marginal likelihood
(the denominator). Fortunately, that still allows us to compute ratios
of posterior densities! We have, \begin{align*}
\frac{p(\theta \mid x)}{p(\theta' \mid x)} = \frac{p(\theta, x)}{p(x)} \frac{p(x)}{p(\theta', x)} = \frac{p(\theta, x)}{p(\theta', x)}.
\end{align*} Now rearrange the detailed balance condition to relate
ratios of transition probabilities to ratios of joint probabilities,
\begin{align*}
\frac{\pi(\theta \mid \theta')}{\pi(\theta' \mid \theta)} = \frac{\pi^\star(\theta)}{\pi^\star(\theta')} 
= \frac{p(\theta \mid x)}{p(\theta' \mid x)} = \frac{p(\theta, x)}{p(\theta', x)}
\end{align*}

\end{frame}

\begin{frame}[t,allowframebreaks]{The Metropolis-Hastings
algorithm}\label{the-metropolis-hastings-algorithm}

To construct such a transition distribution
\(\pi(\theta \mid \theta')\), break it down into two steps.

\begin{enumerate}
\item
  Sample a proposal \(\theta\) from a \emph{proposal distribution}
  \(q(\theta \mid \theta')\),
\item
  Accept the proposal with \emph{acceptance probability}
  \(a(\theta' \to \theta)\). (Otherwise, set \(\theta = \theta'\).)
\end{enumerate}

Thus, \begin{align*}
\pi(\theta \mid \theta') = 
\begin{cases}
q(\theta \mid \theta') \, a(\theta' \to \theta) & \text{if } \theta' \neq \theta \\
\int q(\theta'' \mid \theta') \, (1 - a(\theta' \to \theta'')) \dif \theta'' & \text{if } \theta' = \theta
\end{cases}
\end{align*}

Detailed balance is trivially satisfied when \(\theta = \theta'\). When
\(\theta \neq \theta'\), we need \begin{align*}
\frac{\pi(\theta \mid \theta')}{\pi(\theta' \mid \theta)} = \frac{q(\theta \mid \theta') \, a(\theta' \to \theta)}{q(\theta' \mid \theta) \, a(\theta \to \theta')} = \frac{p(\theta, x)}{p(\theta', x)} \Rightarrow \frac{a(\theta' \to \theta)}{a(\theta \to \theta')} = \underbrace{\frac{p(\theta, x) \, q(\theta' \mid \theta)}{p(\theta', x) \, q(\theta \mid \theta')}}_{\triangleq A(\theta' \to \theta)}
\end{align*}

WLOG, assume $A(\theta' \to \theta) \leq 1$. (If it's not, its
inverse \(A(\theta \to \theta')\) must be.) A simple way to ensure
detailed balance is to set
\(a(\theta' \to \theta) = A(\theta' \to \theta)\) and
\(a(\theta \to \theta') = 1\).

We can succinctly capture both cases with, \begin{align*}
a(\theta' \to \theta) = \min \left\{1, \, A(\theta' \to \theta) \right \} = \min \left\{1, \, \frac{p(\theta, x) \, q(\theta' \mid \theta)}{p(\theta', x) \, q(\theta \mid \theta')} \right \}.
\end{align*}
\end{frame}

\begin{frame}[t,allowframebreaks]{The Metropolis algorithm}\label{the-metropolis-algorithm}

Now consider the special case in which the proposal distribution is
symmetric; i.e.~\(q(\theta \mid \theta') = q(\theta' \mid \theta)\).
Then the proposal densities cancel in the acceptance probability and,
\begin{align*}
a(\theta' \to \theta) = \min \left\{1, \, \frac{p(\theta, x)}{p(\theta', x)} \right \}.
\end{align*} In other words, you accept any proposal that moves
``uphill,'' and only accept ``downhill'' moves with some probability.

This is called the \emph{Metropolis algorithm} and it has close
connections to \emph{simulated annealing}.
\end{frame}

\begin{frame}[t,allowframebreaks]{Smarter Proposals with
Gradients}\label{smarter-proposals-with-gradients}

\begin{itemize}
\item
  Metropolis-Hastings with a symmetric Gaussian proposal behaves (kind
  of) like a random walk.
\item
  Neal (2012) argues that in \(D\) dimensions, random
  walk MH needs \(O(D^2)\) iterations to get an independent sample.
\item
  Can we develop more efficient transition distributions? \textbf{Yes!}
  If we have more information about the log probability.
\item
  For example, suppose that the log probability \(\log p(\theta)\) is
  differentiable. We can use the gradient to make proposals that move
  farther and are more likely to be accepted.
\end{itemize}
\end{frame}

\begin{frame}[t,allowframebreaks]{Metropolis-Adjusted Langevin Algorithm
(MALA)}\label{metropolis-adjusted-langevin-algorithm-mala}

The \emph{Metropolis-Adjusted Langevin Algorithm} uses the gradient of
the log probability to make asymmetric proposals, \begin{align*}
    q(\mbtheta' \mid \mbtheta) &= \cN(\mbtheta + \tau \nabla_{\mbtheta} \log p(\mbtheta, \mbX), \, 2\tau^2 \mbI) 
\end{align*}

\emph{Note:}
\(q(\mbtheta' \mid \mbtheta) \neq q(\mbtheta \mid \mbtheta')\)! To
calculate the acceptance probability, you need the gradient at both
points.

MALA can be motivated as a discrete-time approximation to the
\emph{Langevin} diffusion, a continuous-time stochastic differential
equation for modeling molecular dynamics.

In high dimensions, the extra information provided by the gradient can
lead to much more efficient chains. Neal argues that MALA needs
\(O(D^{4/3})\) computation to produce an independent sample.

But why stop at one gradient step?
\end{frame}

\begin{frame}[t,allowframebreaks]{Hamiltonian Monte Carlo
(HMC)}\label{hamiltonian-monte-carlo-hmc}

\textbf{Idea:} Think of negative log probability as an \emph{energy
landscape}. Now imagine a puck sliding around on this bumpy surface.
Give it random kicks; it will tend to slide downhill toward points of
low potential energy (high probability). Each kick can displace the puck
by a large amount. Done properly, the puck will visit points with
probability proportional to the posterior probability.
\end{frame}

\begin{frame}[t,allowframebreaks]{Notation}\label{notation}

Following Neal (2012), let

\begin{itemize}
\item
  \(\mbq \in \reals^D\) denote the \emph{position}; i.e.~the current
  parameters (previously \(\theta\))
\item
  \(\mbp \in \reals^D\) denote the \emph{momentum}; auxiliary variables
  that we don't care about, but which are necessary for HMC.
\item
  \(\mbz = [\mbq, \mbp]^\top \in \reals^{2D}\) denote the combined
  \emph{state of the system}.
\item
  \(\mbM\) denote the \emph{mass matrix}, another artificial construct.
  Typically, this will be \(m \mbI\)
\item
  \(U(\mbq)\) denote the \emph{potential energy}
\item
  \(K(\mbp) = \frac{1}{2} \mbp^\top \mbM^{-1} \mbp\) denote the
  \emph{kinetic energy}
\end{itemize}
\end{frame}


\begin{frame}[t,allowframebreaks]{Hamiltonian Dynamics}\label{hamiltonian-dynamics}

The \emph{Hamiltonian} is the sum of the potential
\(H(\mbq, \mbp) = U(\mbq) + K(\mbp) = U(\mbq) + \frac{1}{2} \mbp^\top \mbM^{-1} \mbp\).

The partial derivatives determine how the state evolves over time,
\begin{align*}
    \frac{\dif q_d}{\dif t} &= \frac{\partial H}{\partial p_d} = [\mbM^{-1} \mbp]_d \\
    \frac{\dif p_d}{\dif t} &= -\frac{\partial H}{\partial q_d} = -\frac{\partial  U}{\partial q_d}
\end{align*} for \(d=1,\ldots, D\).

Compactly, \begin{align*}
    \frac{\dif \mbz}{\dif t} &= \mbJ \nabla H(z)
    \qquad\text{ where } \qquad
    \mbJ = \begin{bmatrix} \mbzero, \mbI \\ -\mbI, \mbzero
    \end{bmatrix}
\end{align*} 
\end{frame}

    

\begin{frame}[t,allowframebreaks]{Using Hamiltonian Dynamics for Posterior
Inference}\label{using-hamiltonian-dynamics-for-posterior-inference}

Define a joint distribution on positions and momenta as, \begin{align*}
    p(\mbq, \mbp) &\propto \exp \left \{-H(\mbq, \mbp) \right \} \propto 
    \exp \left \{-U(\mbq) - K(\mbp)\right \}.
\end{align*}

Now let \(U(\mbq) = -\log p(\mbtheta = \mbq, \mbX)\) be the
\emph{negative} log joint probability. Then, \begin{align*}
    p(\mbq, \mbp) &= p(\mbtheta = \mbq \mid \mbX) \times p(\mbp)
\end{align*} Samples of \(\mbq\) will be marginally distributed
according to the posterior \(p(\mbtheta = \mbq \mid \mbX)\).

Samples of \(\mbp\) will be marginally distributed
\(p(\mbp) = \frac{\exp \{ -K(\mbp)\}}{\int_{\reals^D} \exp \{-K(\mbp)\} \dif \mbp}\).
These are \emph{auxiliary variables} that we don't really care
about---they're just there to help us construct MH proposals.

We choose \(K(\mbp)\) so \(p(\mbp)\) is convenient; e.g.~if
\(K(\mbp) = \frac{1}{2} \mbp^\top \mbM^{-1} \mbp\) then \begin{align*}
    p(\mbp) &= \cN(\mbp \mid \mbzero, \mbM).
\end{align*}
\end{frame}

\begin{frame}[t,allowframebreaks]{Algorithm}\label{algorithm}

\textbf{Hamiltonian Monte Carlo (HMC)} is Metropolis-Hastings on the
joint distribution of \((\mbq, \mbp)\) with proposals based on
Hamiltonian dynamics.

Starting at point \((\mbq', \mbp')\), sample the proposal distribution:

\begin{enumerate}
\def\labelenumi{\arabic{enumi}.}
\item
  Throw away \(\mbp'\) and sample new momenta from their marginal
  distribution \(\mbp \sim \cN(\mbzero, \mbM)\).
\item
  Approximate Hamiltonian dynamics for \(\Delta t\) time to get to a new
  point \((\mbq, \mbp)\). (Use \(L = \Delta t/\epsilon\) Leapfrog
  integration steps each of size \(\epsilon\).)
\item
  Flip the momentum \(\mbp \leftarrow -\mbp\) to make the proposal
  symmetric.
\end{enumerate}

Then accept the proposed point \((\mbq, \mbp)\) with probability,
\begin{align*}
    a((\mbq', \mbp') \to (\mbq, \mbp)) &= 
    \min \left\{1, \, \frac{\exp\{-H(\mbq, \mbp)\} \, q(\mbq', \mbp' \mid \mbq, \mbp)}{\exp\{-H(\mbq', \mbp')\} q(\mbq, \mbp \mid \mbq', \mbp')} \right \} \\
    &=
    \min \left\{1, \, \frac{\exp\{-H(\mbq, \mbp)\}}{\exp \{-H(\mbq', \mbp')\}} \right \}.
\end{align*} If the Hamiltonian dynamics were simulated exactly, HMC
would always accept. In practice, differences arise from numerical
integration errors.
\end{frame}

\begin{frame}[t,allowframebreaks]{Gibbs Sampling}\label{gibbs-sampling}

Gibbs is a special case of MH with proposals that always accept. Gibbs
sampling updates one ``coordinate'' of \(\theta \in \reals^D\) at a time
by sampling from its conditional distribution. The algorithm is:

\textbf{Gibbs Sampling:}

\textbf{Input:} Initial parameters \(\theta^{(0)}\), observations \(x\)

\begin{itemize}
\item
  \textbf{For} \(t=1,\ldots, T\)

  \begin{itemize}
  \item
    \textbf{For} \(d=1,\ldots, D\)

    \begin{itemize}
    \item
      Sample
      \(\theta_d^{(t)} \sim p(\theta_d \mid \theta_{1}^{(t)}, \ldots, \theta_{d-}^{(t)}, \theta_{d+1}^{(t-1)}, \ldots, \theta_D^{(t-1)}, x)\)
    \end{itemize}
  \end{itemize}
\end{itemize}

\textbf{Return} samples \(\{\theta^{(t)}\}_{t=1}^T\) ```

You can think of Gibbs as cycling through \(D\) Metropolis-Hastings
proposals, one for each coordinate \(d \in 1,\ldots,D\), \begin{align*}
q_d(\theta \mid \theta') = p(\theta_d \mid \theta'_{\neg d}, x) \, \delta_{\theta'_{\neg d}}(\theta_{\neg d}),
\end{align*} where
\(\theta_{\neg d} = (\theta_1, \ldots, \theta_{d-1}, \theta_{d+1}, \ldots, \theta_D)\)
denotes all parameters except \(\theta_d\).

In other words, the proposal distribution \(q_d\) samples \(\theta_d\)
from its conditional distribution and leaves all the other parameters
unchanged.

What is the probability of accepting this proposal? \begin{align*}
a_d(\theta' \to \theta) 
&= \min \left\{ 1, \, \frac{p(\theta, x) q_d(\theta' \mid \theta)}{p(\theta', x) q_d(\theta \mid \theta')} \right\} \\
&= \min \left\{ 1, \, \frac{p(\theta, x) p(\theta_d' \mid \theta_{\neg d}, x) \delta_{\theta_{\neg d}}(\theta'_{\neg d})}{p(\theta', x) p(\theta_d \mid \theta'_{\neg d}, x) \delta_{\theta'_{\neg d}}(\theta_{\neg d})} \right\} \\
&= \min \left\{ 1, \, \frac{p(\theta_{\neg d}, x) p(\theta_d \mid \theta_{\neg d}, x) p(\theta_d' \mid \theta_{\neg d}, x) \delta_{\theta_{\neg d}}(\theta'_{\neg d})}{p(\theta'_{\neg d}, x) p(\theta'_d \mid \theta'_{\neg d}, x) p(\theta_d \mid \theta'_{\neg d}, x) \delta_{\theta'_{\neg d}}(\theta_{\neg d})} \right\} \\
&= \min \left\{1, 1 \right\} = 1
\end{align*} for all \(\theta, \theta'\) that differ only in their
\(d\)-th coordinate.

\textit{The Gibbs proposal is an offer you cannot refuse.}

Of course, if we only update one coordinate, the chain can't be ergodic.
However, if we cycle through coordinates it generally will be.

\textbf{Question:} If Gibbs sampling always accepts, is it strictly
better than other Metropolis-Hastings algorithms? :::
\end{frame}

\begin{frame}[t,allowframebreaks]{Conclusion}\label{conclusion}

This was obviously a whirlwind of an introduction to Bayesian inference!
There's plenty more to be said about Bayesian statistics --- choosing a
prior, subjective vs objective vs empirical Bayesian approaches, the
role of the marginal likelihood in Bayesian model comparison, varieties
of MCMC, and other approaches to approximate Bayesian inference. We'll
dig into some of these topics as the course goes on, but for now, we
have some valuable tools for developing Bayesian modeling and inference
with discrete data!
\end{frame}

    
    
\end{document}
